\documentclass[12pt]{article}
\usepackage[a4paper, margin=2cm]{geometry}

\usepackage{fontspec}
\usepackage{polyglossia}
\setdefaultlanguage{russian}
\setmainlanguage{russian}
\setmainfont{Times New Roman}
\begin{document}

\begin{center}
\textbf{\LARGE ТЕХНИЧЕСКИЙ ОТЧЕТ}
\end{center}

\vspace{1cm}

\noindent Настоящий технический отчет составлен по результатам проверки оборудования (далее --- Оборудование) с целью подтверждения его технического состояния, работоспособности и соответствия установленным требованиям. Пакетная проверка выполнена в условиях эксплуатации/лаборатории с соблюдением установленных регламентов и требований техники безопасности.

\vspace{0.5cm}

\section*{1. Сведения об оборудовании}
\begin{itemize}
    \item Наименование оборудования: \rule{0.75\linewidth}{0.4pt}
    \item Модель/тип: \rule{0.75\linewidth}{0.4pt}
    \item Заводской (серийный) номер: \rule{0.75\linewidth}{0.4pt}
    \item Инвентарный номер (при наличии): \rule{0.75\linewidth}{0.4pt}
    \item Место установки/эксплуатации: \rule{0.75\linewidth}{0.4pt}
    \item Дата проведения проверки: \rule{0.75\linewidth}{0.4pt}
\end{itemize}

\vspace{0.5cm}

\section*{2. Цель и объем проверки}
\noindent Цель проверки: выявление возможных неисправностей, оценка текущего технического состояния, подтверждение корректности функционирования узлов и систем Оборудования, а также формирование заключения о пригодности к дальнейшей эксплуатации.

\noindent Объем проверки включает:
\begin{itemize}
    \item визуальный осмотр и контроль целостности узлов;
    \item проверку внешних признаков повреждений, загрязнений, следов перегрева/коррозии;
    \item функциональные проверки в заданных режимах;
    \item контроль параметров в соответствии с эксплуатационной документацией (при наличии);
    \item оформление результатов и выводов.
\end{itemize}

\vspace{0.5cm}

\section*{3. Используемые средства и условия проведения}
\noindent Для выполнения проверки применялись измерительные приборы и вспомогательное оборудование (перечень приведен ниже либо по внутреннему реестру средств измерений).

\begin{itemize}
    \item Наименование измерительных средств: \rule{0.9\linewidth}{0.4pt}
    \item Тип/модель: \rule{0.9\linewidth}{0.4pt}
    \item Серийный номер: \rule{0.9\linewidth}{0.4pt}
    \item Дата поверки/калибровки (при наличии): \rule{0.7\linewidth}{0.4pt}
    \item Условия проведения: температура \rule{0.4\linewidth}{0.4pt}, относительная влажность \rule{0.4\linewidth}{0.4pt}, прочие условия \rule{0.5\linewidth}{0.4pt}
\end{itemize}

\vspace{0.5cm}

\section*{4. Ход проверки и результаты измерений}
\noindent Проверка выполнена в следующей последовательности:

\subsection*{4.1. Визуальный осмотр}
\noindent В ходе осмотра установлено:
\begin{itemize}
    \item внешние поверхности: \rule{0.85\linewidth}{0.4pt}
    \item элементы крепления и соединения: \rule{0.85\linewidth}{0.4pt}
    \item наличие/отсутствие следов перегрева, механических повреждений, коррозии: \rule{0.85\linewidth}{0.4pt}
    \item состояние кабельных линий/разъемов (при наличии): \rule{0.85\linewidth}{0.4pt}
\end{itemize}

\subsection*{4.2. Функциональные проверки}
\noindent Выполнены проверки в соответствии с эксплуатационной документацией и/или установленной методикой, включая запуск, контроль режимов работы, оценку стабильности параметров и отсутствие аварийных индикаций. Результаты:
\begin{itemize}
    \item Режим работы 1: \rule{0.85\linewidth}{0.4pt}
    \item Режим работы 2: \rule{0.85\linewidth}{0.4pt}
    \item Дополнительные проверки (если применимо): \rule{0.85\linewidth}{0.4pt}
\end{itemize}

\subsection*{4.3. Измеренные параметры (при наличии)}
\noindent При необходимости выполнялись измерения электрических/механических/технических параметров. Значения представлены в обобщенном виде:
\begin{itemize}
    \item Параметр 1: фактическое значение \rule{0.75\linewidth}{0.4pt}, отклонение/оценка \rule{0.75\linewidth}{0.4pt}
    \item Параметр 2: фактическое значение \rule{0.75\linewidth}{0.4pt}, отклонение/оценка \rule{0.75\linewidth}{0.4pt}
    \item Параметр 3: фактическое значение \rule{0.75\linewidth}{0.4pt}, отклонение/оценка \rule{0.75\linewidth}{0.4pt}
\end{itemize}

\vspace{0.5cm}

\section*{5. Анализ выявленных несоответствий (при наличии)}
\noindent Несоответствия/замечания (если выявлены) и их возможные причины:
\begin{itemize}
    \item \rule{0.95\linewidth}{0.4pt} (возможная причина: \rule{0.6\linewidth}{0.4pt})
    \item \rule{0.95\linewidth}{0.4pt} (возможная причина: \rule{0.6\linewidth}{0.4pt})
\end{itemize}

\noindent Предусмотренные действия (при необходимости):
\begin{itemize}
    \item \rule{0.95\linewidth}{0.4pt}
    \item \rule{0.95\linewidth}{0.4pt}
\end{itemize}

\vspace{0.5cm}

\section*{6. Заключение}
\noindent По результатам проведенной технической проверки установлено, что:
\begin{itemize}
    \item Оборудование: \rule{0.95\linewidth}{0.4pt}
    \item Функционирование в проверенных режимах: \rule{0.95\linewidth}{0.4pt}
    \item Выявленные отклонения (при наличии): \rule{0.95\linewidth}{0.4pt}
\end{itemize}

\noindent Оборудование рекомендуется к эксплуатации/к дальнейшему использованию при условии выполнения мероприятий, указанных в разделе 5 (при наличии). В случае наличия замечаний требуется проведение повторной проверки после выполнения корректирующих действий.

\vspace{0.5cm}

\section*{7. Применимые документы}
\noindent При формировании настоящего отчета использовались:
\begin{itemize}
    \item эксплуатационная документация изготовителя Оборудования;
    \item внутренние регламенты организации по техническому обслуживанию и проверке оборудования;
    \item иные нормативные/методические документы (при наличии): \rule{0.9\linewidth}{0.4pt}
\end{itemize}

\vfill

\noindent
Дата: \rule{5cm}{0.4pt} \hfill Подпись: \rule{5cm}{0.4pt}

\end{document}